% =================
% LEGENDA DOS CÓDIGOS — sistema visual R1076
% =================
\chapter*{Legenda dos Códigos}
\addcontentsline{toc}{chapter}{Legenda dos Códigos}
\markboth{Legenda dos Códigos}{}

\noindent Cada fragmento pertence a uma família. A cor e a capitular de abertura
indicam essa família antes de o leitor começar a ler.

\vspace{1.2em}
\begin{description}
  \item[\textcolor{mem}{\large\bfseries MEM}] \textbf{Memória.} A voz de Joaquim: sertão, fome, colônia, o diário. Cor: cobre queimado.
  \item[\textcolor{doc}{\large\bfseries DOC}] \textbf{Documento.} Laudos, prontuários, cartas, registros e provas. Cor: teal profundo.
  \item[\textcolor{luc}{\large\bfseries LUC}] \textbf{Lucidez.} A investigação de Lúcia Mendes em 2026. Cor: violeta.
  \item[\textcolor{cont}{\large\bfseries CONT}] \textbf{Contaminação.} A segunda pessoa, o leitor, o ciclo. Cor: carmim.
\end{description}

\vspace{1.2em}
\noindent\textit{O Epílogo usa o âmbar do arquivo; a capitular de cada fragmento é a inicial da
família (\textbf{M}emória, \textbf{D}ocumento, \textbf{L}ucidez,
\textbf{C}ontaminação), e o filete abaixo do título usa a mesma cor.}

\vspace{1em}
\noindent\textit{Contraste medido sobre o papel sépia: MEM 6,76:1; DOC 6,76:1;
LUC 9,23:1; CONT 7,18:1. Todas as cores de leitura passam o mínimo de 4,5:1.}
